\section{Signal and Execution Model}
\label{sec:signal-model}
Let $x[n]$ be a real signal sampled at $f_s$\,Hz. A \emph{sample call}
advances the engine by one sample; a host callback, or \emph{block}, renders
$D$ such samples. The two are distinct measurement boundaries. The hop $H$ is
the interval between successive frame endpoints and is unrelated to the block
size. The transform length $N\ge4$ is a power of two, and $H\ge1$ is an
integer. With continuously captured input and a fixed hop, frame $r$ ends at
$t_r=rH$ and analyzes the window
\begin{equation}
 u_r[m]=x[t_r-N+1+m]w[m],\qquad 0\le m<N,
\end{equation}
where $w$ is the analysis window. Its unnormalized forward transform is
\begin{equation}
 X_r[k]=\sum_{m=0}^{N-1}u_r[m]e^{-j2\pi km/N},
 \label{eq:dft}
\end{equation}
with bin $k$ at frequency $kf_s/N$. Deferring computation changes when $X_r$
becomes available, never which samples enter this sum. We define
\emph{sample-index age} as the distance from a specified frame timestamp to
publication, converted to seconds using $f_s$. It excludes wall-clock
computation, device buffering, and display consumption. In particular, a
transform computed within its endpoint's sample call has zero
endpoint-to-publication age but nonzero execution time.

For amplitude calibration, the finite-window coherent gain is
$G_c=N^{-1}\sum_m w[m]$. An isolated, bin-centered real sinusoid at an
interior bin has amplitude $2|X_r[k]|/(NG_c)$, provided the window's
positive- and negative-frequency contributions do not overlap materially; DC
and Nyquist omit the factor of two. This is amplitude calibration, not
power-spectral-density normalization, which depends on the window's energy
instead. The plugin caches window samples and applies tabulated
coherent-gain factors, which should not be assumed to equal the exact
finite-length sum for every window and length. The numerical audits in
Section~\ref{sec:evaluation} compare each path with independent references
and therefore verify computation, not display calibration.

Real-input packing uses $M=N/2$ complex samples, and the analyzer retains the
$K=M+1$ nonnegative-frequency bins, including DC and Nyquist. The standard
radix-2 traversal of the packed transform requires
$B=B_r(N)=(N/4)\log_2(N/2)$ butterflies. Appendix~\ref{sec:legacy-transform}
gives the butterfly equations, the arithmetic-order proof, the resumable
state transition, and real-spectrum reconstruction.
